\section{Related Work}

\paragraph{Cough and respiratory audio datasets.}
COUGHVID provides a large-scale crowdsourced cough corpus with expert annotation for a subset of recordings \cite{orlandic2021coughvid}. Coswara extends cough-only screening with breathing, phonation, speech, demographics, symptoms, and SARS-CoV-2 status \cite{bhattacharya2023coswara}. Cambridge COVID-19 Sounds and the DiCOVA challenges highlighted a central methodological issue: cough models can perform strongly under narrow splits while degrading under realistic external validation \cite{han2022sounds}. ICBHI 2017 is not a cough dataset, but it remains useful for respiratory-sound transfer and robustness analysis.

\paragraph{Cough-based disease screening.}
Earlier cough systems used MFCC, spectral descriptors, bag-of-words, SVMs, random forests, and shallow neural networks \cite{pavel2023bow,pahar2021cough}. Later systems adopted CNNs, LSTMs, spectrogram image classifiers, and ensemble learning \cite{hussain2024cough2covid}. Across this literature, AUROC, sensitivity, specificity, F1, and subgroup performance are more informative than accuracy alone because cough datasets are often imbalanced and clinically asymmetric. Recent high reported accuracies should be interpreted carefully unless subject-level splits, external datasets, and threshold selection policies are explicit.

\paragraph{Audio foundation models.}
PANNs demonstrated that CNNs pretrained on AudioSet transfer well across audio tasks \cite{kong2020panns}. AST adapted the Vision Transformer to spectrogram patches and showed the strength of attention-only audio classification \cite{gong2021ast}. PaSST improved transformer training efficiency through Patchout \cite{koutini2022passt}. HTS-AT introduced hierarchical token-semantic structure for sound classification and detection \cite{chen2022htsat}. BEATs used acoustic tokenizers for iterative self-supervised audio pretraining and achieved strong AudioSet and downstream performance \cite{chen2023beats}. These models are strong teachers for cough screening, but they are not ideal endpoint models for low-power devices.

\paragraph{Efficient audio students.}
MobileNetV3 and EfficientNet are established compact vision backbones that transfer naturally to spectrogram inputs \cite{howard2019mobilenetv3,tan2019efficientnet}. BC-ResNet was designed for efficient keyword spotting and is attractive for tiny audio classifiers \cite{kim2021bcresnet}. ECAPA-TDNN is stronger for speaker embeddings than disease labels, but its channel attention and temporal pooling make it a useful compact baseline for cough representations. The main empirical question is not whether small models run faster, but whether they retain teacher robustness under dataset shift.

\paragraph{Knowledge distillation.}
Response distillation transfers softened teacher logits \cite{hinton2015distilling}. Feature distillation aligns hidden representations \cite{romero2015fitnets}. Attention transfer encourages similar saliency or activation distributions \cite{zagoruyko2017attention}. Relation and contrastive distillation preserve pairwise or class-structured geometry \cite{tian2020crd,khosla2020supcon}. Recent acoustic monitoring work also shows KD can support efficient audio detection pipelines outside standard vision/NLP settings \cite{priebe2024speechkd}. \textsc{CoughKD} combines these ideas because cough labels are noisy and sparse: response KD captures class similarity, feature KD transfers acoustic abstractions, attention KD regularizes time-frequency focus, and relation KD stabilizes representation geometry across domains.
